\section{Introduction}

Frequently, tasks in machine learning can be expressed as the problem of
optimizing an objective function $f(\theta)$ defined over some domain
$\theta\in\Theta$. The goal in this case is to find the minimizer
$\theta^* =
\argmin_{\theta\in\Theta}f(\theta)$. While any method capable of minimizing
this objective function can be applied, the standard approach for
differentiable functions is some form of gradient descent, resulting in a
sequence of updates
\begin{align*}
  \theta_{t+1} = \theta_t - \alpha_t \nabla f(\theta_t)
  \,.
\end{align*}
The performance of vanilla gradient descent, however, is hampered by the fact
that it \emph{only} makes use of gradients and ignores second-order
information. Classical optimization techniques correct this behavior by
rescaling the gradient step using curvature information, typically via the
Hessian matrix of
second-order partial derivatives---although other choices such as the
generalized Gauss-Newton matrix or Fisher information matrix are possible.

Much of the modern work in optimization is based around designing update rules
tailored to specific classes of problems, with the types of problems of
interest differing between different research communities.  For example, in the
deep learning community we have seen a proliferation of optimization methods
specialized for high-dimensional, non-convex optimization problems. These
include momentum \citep{nesterov:1983,tseng:1998}, Rprop \citep{riedmiller:1992},
Adagrad \citep{duchi:2011}, RMSprop \citep{tieleman:2012}, and ADAM \citep{kingma:2015}. More
focused methods can also be applied when more structure of the optimization
problem is known
\citep{martens:2015}. In contrast, communities who focus on sparsity tend to
favor very different approaches~\citep{donoho:2006,bach:2012}. This is even
more the case for combinatorial optimization for which relaxations are often
the norm \citep{nemhauser:1988}.

\begin{wrapfigure}[18]{r}{0.5\textwidth}
\centering
\includegraphics[width=\linewidth]{figures/optimizer-optimizee}
\caption{The optimizer (left) is provided with performance of the optimizee
(right) and proposes updates to increase the optimizee's performance.
\citep[photos:][]{neuron1,neuron2}}
\label{fig:optimizee}
\end{wrapfigure}

This industry of optimizer design allows different communities to create
optimization methods which exploit structure in their problems of interest at
the expense of potentially poor performance on problems outside of that scope.
Moreover the \emph{No Free Lunch Theorems for Optimization}
\citep{wolpert:1997} show that in the setting of combinatorial optimization, no
algorithm is able to do better than a random strategy in expectation.  This
suggests that specialization to a subclass of problems is in fact the
\emph{only} way that improved performance can be achieved in general.

% In this setting \citeauthor{wolpert:1997} consider
% combinatorial optimization problems and show that in expectation (over the set
% of all optimization problems) no algorithm is able to do better than a random
% strategy. It is only by considering a suitable subset and tailoring the
% algorithm to that subset is improved performance attainable---in particular
% this new algorithm \emph{must exhibit worse performance} on problems outside
% the class.


In this work we take a different tack and instead propose to replace hand-designed update rules with a learned update rule, which we call the optimizer
$g$, specified by its own set of parameters $\phi$. This results in updates to
the optimizee $f$ of the form
\begin{equation}
  \theta_{t+1} = \theta_t + g_t(\nabla f(\theta_t), \phi)
  \,.
  \label{eq:meta-update}
\end{equation}
A high level view of this process is shown in Figure~\ref{fig:optimizee}.  In
what follows we will explicitly model the update rule $g$ using a recurrent
neural network (RNN) which maintains its own state and hence dynamically
updates as a function of its iterates.

% Given a fixed set of
% parameters $\phi$ the learning rule $g_t$ will then act like a form of gradient
% descent and we will refer to this process as \emph{learning by gradient
% descent}. Here the base learning procedure has the same semantics as are
% typically used in machine learning, where $f$ may represent an average loss
% over some supervised dataset.

% The process of learning a learning rule has a long history,
% generally going under the name of \emph{learning to learn} or \emph{meta-learning}
% \citep{thrun:1998}.
% Training $\phi$ results in a gradient-based optimization procedure, so we call
% this \emph{learning to learn by gradient descent}, and given that we also train
% $\phi$ using gradient descent we call the full procedure \emph{learning to
% learn by gradient descent by gradient descent}. This formalism allows us to
% train the learning mechanism restricted to a suitable set of example functions
% $f$ and hence tailor the rule to this particular subset.


\subsection{Transfer learning and generalization}

The goal of this work is to develop a procedure for constructing a learning
algorithm which performs well on a particular class of optimization problems.
Casting algorithm design as a learning problem allows us to specify the class
of problems we are interested in through example problem instances.  This is in
contrast to the ordinary approach of characterizing properties of interesting
problems analytically and using these analytical insights to design learning
algorithms by hand.

It is informative to consider the meaning of \emph{generalization} in this
framework.  In ordinary statistical learning we have a particular function of
interest, whose behavior is constrained through a data set of example function
evaluations.  In choosing a model we specify a set of inductive biases about
how we think the function of interest should behave at points we have not
observed, and generalization corresponds to the capacity to make predictions
about the behavior of the target function at novel points.
In our setting the examples are themselves \emph{problem instances}, which
means generalization corresponds to the ability to transfer knowledge between
different problems.  This reuse of problem structure is commonly known as
\emph{transfer learning}, and is often treated as a subject in its own right.
However, by taking a meta-learning perspective, we can cast the problem of
transfer learning as one of generalization, which is much better studied in the
machine learning community.

One of the great success stories of deep-learning is that we can rely on the
ability of deep networks to generalize to new examples by learning interesting
sub-structures. In this work we aim to leverage this generalization power, but
also to lift it from simple supervised learning to the more general setting of
optimization.

\subsection{A brief history and related work}

The idea of using \emph{learning to learn} or \emph{meta-learning} to acquire
knowledge or inductive biases has a long history \citep{thrun:1998}. More
recently, \cite{lake:2016} have argued forcefully for its importance as a
building block in artificial intelligence. Similarly, \cite{santoro:2016}
frame multi-task learning as generalization, however unlike our approach they
directly train a base learner rather than a training algorithm.  In general
these ideas involve learning which occurs at two different time scales: rapid
learning within tasks and more gradual, \emph{meta} learning across many
different tasks.

Perhaps the most general approach to meta-learning is that of
\cite{schmidhuber:1992, schmidhuber:1993}---building on work from
\citep{schmidhuber:1987}---which considers networks that are able to modify
their own weights. Such a system is differentiable end-to-end, allowing both
the network and the learning algorithm to be trained jointly by gradient
descent with few restrictions.  However this generality comes at the expense of
making the learning rules very difficult to train.  Alternatively, the work of
\cite{schmidhuber:1997} uses the Success Story Algorithm to modify its search
strategy rather than gradient descent; a similar approach has been recently
taken in \cite{daniel:2016} which uses reinforcement learning to train a
controller for selecting step-sizes.

\cite{bengio:1990, bengio:1995} propose to learn updates which avoid
back-propagation by using simple parametric rules. In relation to the focus of
this paper the work of \citeauthor{bengio:1990} could be characterized as
\emph{learning to learn \textbf{without} gradient descent by gradient descent}.
The work of \cite{runarsson:2000} builds upon this work by replacing the simple
rule with a neural network.

\cite{cotter:1990}, and later \cite{younger:1999}, also show fixed-weight
recurrent neural networks can exhibit dynamic behavior without need to modify
their network weights. Similarly this has been shown in a filtering context
\citep[e.g.][]{feldkamp:1998}, which is directly related to simple
multi-timescale optimizers \citep{sutton:1992, schraudolph:1999}.

Finally, the work of \cite{younger:2001} and \cite{hochreiter:2001} connects
these different threads of research by allowing for the output of
backpropagation from one network to feed into an additional \emph{learning}
network, with both networks trained jointly. Our approach to meta-learning
builds on this work by modifying the network architecture of the optimizer in
order to scale this approach to larger neural-network optimization problems.

